Un ensemble de circuits électriques et électroniques comportent des
condensateurs et des bobines. Le comportement de ces circuits diffère selon
l'effet imposé par ces composants. Le but de cet exercice est d'étudier un
circuit RLC série dans différents cas.

\begin{minipage}{0.70\linewidth}
    On réalise le montage expérimental représenté sur la
    Fig.~\ref{fig:2bac_svt_national_2017_normale_phys_ex2_1} qui comporte~:
    \begin{itemize}
        \item un générateur idéal de tension de force électromotrice
              ${E=\qty{6}{\volt}}$~;
        \item un condensateur de capacité $C$~;
        \item un conducteur ohmique de résistance $R$~;
        \item une bobine $b$ d'inductance $L$ et de résistance $r$~;
        \item un interrupteur $\mathrm{K}$.
    \end{itemize}
\end{minipage}
\hfill
\begin{minipage}{0.28\linewidth}
    \flushright
    \begin{figure}[H]
    \begin{circuitikz}
        \newdimen\d \d=3cm
        \draw (0,0) node[cute spdt mid,rotate=90] (Sw) {};
        \node[above,font=\ttfamily\small] at (Sw.out 1) {1};
        \node[above,font=\ttfamily\small] at (Sw.out 2) {2};
        \node[above=9mm] at (Sw.in) {K};
        \draw (Sw.out 1)
             to[R,l_=$R$,voltage shift=2.5] ++(-0.7\d,0)
             to[vsource,v_<=$E$,i<^=$i$,voltage shift=0.5] ++(0,-\d)
             to[short] ++(1.4\d,0) coordinate(M)
             to[L,l={$(L,r)$},mirror,voltage shift=3.5,name=C] ++(0,\d)
             to[short] (Sw.out 2);
        \draw (Sw.in) to[C,l_=$C$] (Sw.in |- M);
    \end{circuitikz}
    \caption{}
    \label{fig:2bac_svt_national_2017_normale_phys_ex2_1}
    \end{figure}
\end{minipage}

\begin{enumerate}
    \item On place l'interrupteur dans la position~(1), le condensateur se
          charge totalement. Sa charge maximale est
          $Q_{\max }=\qty{1,32e-4}{\coulomb}$.
          
          
          Calculer la valeur de l'énergie électrique maximale $E_{e,\max}$
          emmagasinée dans le condensateur.
    \item On réalise trois expériences en utilisant trois bobines différentes
          $b_1$, $b_2$ et $b_3$ dont les caractéristiques sont~:
          \[
              b_1\left(L_1=\qty{260}{\mH}~; r_1=0\right), \quad
              b_2\left(L_2=\qty{115}{\mH}~; r_2=0\right) \quad \text{et} \quad
              b_3\left(L_3~; r_3=\qty{10}{\ohm}\right)
          \]

          Dans chaque expérience, on charge totalement le condensateur et on le
          décharge dans l'une des trois bobines.

          Les courbes de la
          Fig.~\ref{fig:2bac_svt_national_2017_normale_phys_ex2_2} représentent
          les variations de la tension $u_C(t)$ aux bornes du condensateur.

          \begin{figure}[H]
              \centering
              \includegraphics[width=0.8\linewidth]{2bac_svt_national_2017_normale_phys_ex2_2}
              \caption{}
              \label{fig:2bac_svt_national_2017_normale_phys_ex2_2}
          \end{figure}

          \begin{enumerate}
              \item Nommer les régimes d'oscillations mis en évidence par les
                    courbes~(a) et~(c).
              \item En comparant les périodes des oscillations électriques,
                    montrer que la courbe~(a) correspond à la bobine $b_2$.
              \item Vérifier que $C \simeq \qty{2,2e-5}{\farad}$.
          \end{enumerate}
    \item On considère le cas de la décharge du condensateur à travers la bobine
          $b_2\left(L_2=\qty{115}{\mH}~; r_2=0\right)$.

          Dans ce cas le circuit LC est idéal.
          \begin{enumerate}
              \item Établir l'équation différentielle vérifiée par la tension
                    $u_C(t)$.
              \item La solution de l'équation différentielle s'écrit~:
                    $u_c(t)=U_{\mathrm{C,\max}} \cdot
                    \cos\left(\frac{2 \pi}{T_0}\cdot t+\varphi\right)$.
                    \begin{enumerate}
                        \item Écrire l'expression numérique de la tension
                              $u_c(t)$.
                        \item Calculer l'énergie totale du circuit LC sachant
                              qu'elle se conserve.
                    \end{enumerate}
          \end{enumerate}
    \item On considère le cas de la décharge du condensateur à travers la bobine
          $b_3\left(L_3~; r_3=\qty{10}{\ohm}\right)$.

          Pour entretenir les oscillations électriques, on ajoute au circuit un
          générateur qui délivre une tension proportionnelle à l'intensité du
          courant $u_g=k\cdot i(t)$ où $k$ est une constante positive. On obtient
          des oscillations électriques sinusoïdales de période
          $T=\qty{10}{\ms}$.
          \begin{enumerate}
              \item Déterminer la valeur de $k$.
			  \item En déduire la valeur de $L_3$.
          \end{enumerate}
\end{enumerate}

